\section{Langue, corpus et modèles embarqués}

\subsection{Variété cible et méthode linguistique}

KORA cible le \textbf{fulfulde de l’Adamaoua}, code ISO 639-3
\code{fub}, en écriture latine. Les caractères \texttt{ɓ}, \texttt{ɗ},
\texttt{ƴ}, \texttt{ŋ} et \texttt{ɲ} sont conservés en UTF-8. Les corpus
étiquetés \code{ff}, \code{fuv}, \code{fuf} ou \code{ffm} ne sont pas
silencieusement assimilés à cette variété camerounaise.

\begin{figure}[H]
\centering
\begin{minipage}[c]{0.47\linewidth}
  \includegraphics[width=\linewidth]{assets/validation.png}
\end{minipage}
\hfill
\begin{minipage}[c]{0.49\linewidth}
  \textbf{Processus appliqué à une entrée critique}
  \begin{enumerate}
    \item identifier le concept financier en FR et EN;
    \item rechercher une forme camerounaise \code{fub};
    \item recouper l’orthographe et le contexte;
    \item distinguer attestation, emprunt et composition;
    \item vérifier que FR et EN donnent la même sortie;
    \item tester les valeurs, négations et actions;
    \item enregistrer source, statut et décision.
  \end{enumerate}
\end{minipage}
\caption{Validation humaine, traçabilité et contrôles automatisés sont
complémentaires.}
\end{figure}

\textbf{\LinguisticValidator} assure la responsabilité de la validation
linguistique du projet. Les formulations historiques relues sont distinguées
des propositions compositionnelles encore à valider. Cette séparation évite
de présenter une génération de corpus comme une norme linguistique.

\subsection{Sources et propriété intellectuelle}

\begin{longtable}{@{}L{43mm}L{29mm}L{43mm}L{37mm}@{}}
\toprule
\textbf{Ressource} & \textbf{Licence} & \textbf{Usage KORA} &
\textbf{Redistribution}\\
\midrule
\endfirsthead
\toprule
\textbf{Ressource} & \textbf{Licence} & \textbf{Usage KORA} &
\textbf{Redistribution}\\
\midrule
\endhead
Common Voice 27 \code{fub} &
CC0 &
1\,067 phrases pour orthographe et modèle de fluidité &
Pas de corpus brut dans l’extension.\\
Tourneux \& Konaï, vocabulaire mobile &
CC BY 4.0 &
Terminologie camerounaise des communications et transferts &
Attribution conservée.\\
Ressources SIL \code{fub} &
Conditions propres &
Consultation et contrôle, sans extraction automatisée &
Non redistribuées.\\
Deux inventaires financiers fournis &
Licence non déclarée &
Inventaire de concepts et définitions françaises &
Fichiers conservés dans le dépôt; pas présentés comme corpus parallèle.\\
PanLex \code{fub} &
Ambiguïté de licence / restriction non commerciale &
Évalué, non intégré &
Non redistribué.\\
Google Translate vers \code{ff} &
Service externe de comparaison &
Cinq réponses publiques comparées; aucune importation automatique &
Résultats archivés séparément.\\
Code KORA &
MIT &
Extension, moteur, modèles et tests &
Distribué sous MIT.\\
\bottomrule
\end{longtable}

\subsection{Inventaires financiers et statuts}

Les deux volumes fournis totalisent 2\,859 lignes et 512 libellés, ramenés à
504 concepts après déduplication de huit chevauchements. Ils décrivent les
concepts en français, mais ne constituent pas une vérité parallèle
FR/EN$\rightarrow$\code{fub}. La matrice synthétique finale du volume 2 est
exclue : elle contient notamment des associations contradictoires entre
réussites et causes d’échec.

\begin{tabularx}{\linewidth}{@{}L{47mm}C{22mm}Y@{}}
\toprule
\textbf{Statut} & \textbf{Nombre} & \textbf{Interprétation}\\
\midrule
\code{reviewed-existing} & 20 &
Formulation déjà présente dans la mémoire KORA relue.\\
\code{compositional-candidate} & 340 &
Composition transparente mais encore à confirmer par revue native.\\
\code{local-loan-candidate} & 31 &
Emprunt technique cohérent avec l’usage local, à confirmer.\\
\code{preserved-code} & 113 &
Code ou identifiant qui reste strictement identique.\\
\midrule
\textbf{Total} & \textbf{504} & 391 clés françaises et 390 clés anglaises.\\
\bottomrule
\end{tabularx}

Les candidats élargissent la reconnaissance de l’interface et alimentent le
modèle à poids réduit, mais leur statut demeure visible. Ils ne remplacent pas
une attestation indépendante.

\subsection{Corrections morphologiques intégrées}

La revue du 14 novembre 2026 a corrigé des faux-sens et harmonisé les noms
d’action. Les concepts emploient préférentiellement la forme nominale
\texttt{-ugol}; une commande peut garder la forme verbale Adamawa en
\texttt{-ugo}.

\begin{tabularx}{\linewidth}{@{}L{38mm}L{49mm}Y@{}}
\toprule
\textbf{Concept} & \textbf{Forme retenue} & \textbf{Décision}\\
\midrule
Compte bloqué & \texttt{Konte uddaaɗo} &
Participe passé adjectival, pas verbe conjugué.\\
Date de valeur & \texttt{Ñalngu kiisaangu} &
Écarte le faux-sens «\,ancien/usé\,».\\
Découvert autorisé & \texttt{Kaalisi ɓeydaaɗi yamiraaɗi} &
Accord avec la classe de \texttt{kaalisi}.\\
Paiement récurrent & \texttt{Yoɓugol lortotoongol} &
Écarte le sens «\,défaire/résilier\,».\\
Paiement sans contact & \texttt{Yoɓugol tawa meemaaka} &
Voyelle longue de \texttt{meemugo}.\\
Remboursement de prêt & \texttt{Yoɓugol nyamaande} &
Sens dette/prêt explicite.\\
Remboursement marchand & \texttt{Ruttinangol kaalisi} &
Restitution d’argent distincte du prêt.\\
Rétrofacturation & \texttt{Ruttinangol kaalisi yoɓaaɗi} &
Nom d’action et objet financier explicites.\\
\bottomrule
\end{tabularx}

\subsection{Pourquoi plusieurs couches plutôt qu’un traducteur génératif}

Deux modèles téléchargeables ont été évalués sur cinq actions critiques :
Flagship \code{francais-fufulde} et DS4H \code{linguo\_mt\_en\_fub}. Les
sorties observées comportaient répétitions, omissions ou contresens; leurs
empreintes dépassaient aussi la limite de 10 Mio. NLLB-200 et M2M-100 sont
beaucoup plus volumineux et ne ciblent pas précisément \code{fub}.

KORA sépare donc le \textbf{sens} de la \textbf{fluidité}. Le moteur
déterministe choisit un concept et protège ses valeurs. Les modèles
linguistiques embarqués ne voient ensuite que des formulations \code{fub}
déjà sûres, qu’ils peuvent classer mais jamais autoriser seuls.

\subsection{Architecture du modèle neuronal}

\begin{figure}[H]
\centering
\resizebox{\linewidth}{!}{%
\begin{tikzpicture}[
  node distance=6mm,
  b/.style={draw=Line,fill=Paper,rounded corners=1.5mm,text width=30mm,
    minimum height=12mm,align=center,font=\small},
  m/.style={b,draw=Teal,fill=Mint},
  arrow/.style={-{Latex[length=2mm]},thick,color=TealDark}
]
\node[b] (ctx) {Trois mots précédents};
\node[m,right=of ctx] (emb) {Embeddings\\24 dimensions};
\node[m,right=of emb] (hid) {Couche tanh\\96 neurones};
\node[m,right=of hid] (out) {Distribution du mot suivant\\vocabulaire 1\,472};
\node[b,below=9mm of hid] (ng) {Trigrammes, bigrammes et unigrammes};
\node[m,right=of out] (rank) {Prédiction et score de fluidité};
\draw[arrow] (ctx)--(emb);
\draw[arrow] (emb)--(hid);
\draw[arrow] (hid)--(out);
\draw[arrow] (out)--(rank);
\draw[arrow] (ng)-|(rank);
\end{tikzpicture}}
\caption{Réseau KORA quantifié et repli n-gramme déterministe.}
\end{figure}

\begin{tabularx}{\linewidth}{@{}L{50mm}Y@{}}
\toprule
\textbf{Caractéristique} & \textbf{Valeur reproductible}\\
\midrule
Architecture & Contexte 3, embedding 24, tanh 96\\
Paramètres & 185\,120\\
Quantification & Poids int8; biais float32\\
Vocabulaire neuronal & 1\,472 unités\\
Données d’entraînement &
1\,067 phrases Common Voice; 9\,368 tokens; 262 formulations KORA;
386 formulations terminologiques à poids réduit\\
Repli n-gramme & 1\,199 mots fréquents; 1\,561 contextes\\
Validation tenue à part & 107 phrases Common Voice\\
Résultat int8 & top-1 14,8\,\%; top-5 35,0\,\%\\
Entraînement & Graine 2026; meilleur état époque 2; arrêt après 8/40\\
Taille générée & Environ 262 Kio avec vocabulaire et métadonnées\\
\bottomrule
\end{tabularx}

La configuration retenue maximise la top-5 parmi les variantes évaluées,
métrique cohérente avec un rôle de classement de plusieurs candidats. Le
score du mot suivant fusionne 85\,\% de réseau et 15\,\% de n-gramme; le score
de phrase combine ensuite 80\,\% de score neuronal et 20\,\% de score
n-gramme. Un cache borné garde les 1\,000 derniers scores.

\begin{warningbox}[Ce que le modèle ne fait pas]
Le modèle de langue n’est pas un traducteur FR/EN autonome. Il ne connaît pas
les règles de conservation financière et ne doit pas être appelé seul pour
produire un message bancaire. La fonction \code{translateText} orchestre la
mémoire, le correcteur, les gabarits, le modèle et le verrou d’intégrité.
\end{warningbox}

\subsection{Preuve d’usage dans l’extension}

\begin{enumerate}
  \item \code{translateText} obtient un candidat dont le sens et les valeurs
        sont déjà contrôlés;
  \item \code{scoreFubFluency} appelle \code{scoreNeuralFub};
  \item l’inférence JavaScript décode les poids int8 et calcule les
        probabilités;
  \item le n-gramme complète les séquences peu vues;
  \item \code{rankFubCandidates} trie uniquement ces candidats sûrs;
  \item le contrôle d’intégrité peut encore rejeter le résultat.
\end{enumerate}

Le fichier \code{src/fub-neural-model.js} est regroupé dans
\code{content.js} au build. Il n’est ni téléchargé au premier démarrage, ni
simulé par une métadonnée. \code{npm run model:demo} et les tests unitaires
appellent les mêmes poids que le bundle Chrome.
